\documentclass[10pt,letterpaper]{article}
\usepackage[margin=0.28in]{geometry}
\usepackage{graphicx}
\usepackage{caption}
\usepackage{amsmath}
\pagestyle{empty}
\setlength{\parindent}{0pt}

\begin{document}
{\Large\bfseries 16-662 Robot Autonomy Homework \#1\par}
\noindent\hfill Name: Aaron Chou, Andrew ID: aaroncho
\vspace{0.1in}

\large \textbf{2 - PID Control}

\begin{minipage}[t]{0.495\textwidth}
  \includegraphics[width=\linewidth]{../../HW1/img/PID_static.png}
  \captionof{figure}{Control with a Fixed Wall}
  \label{fig:fixed-wall}
\end{minipage}\hfill
\begin{minipage}[t]{0.495\textwidth}
  \includegraphics[width=\linewidth]{../../HW1/img/PID_ocillating.png}
  \captionof{figure}{Control with a Moving Wall}
  \label{fig:moving-wall}
\end{minipage}

\vspace{0.1in}

\begin{figure}[htbp]
  \centering
  \includegraphics[width=0.5\linewidth]{../../HW1/img/PID_max_amp.png}
  \caption{Whiteboard near maximum amplitude}\par
  \label{fig:max-amp}
\end{figure}
\vspace{0.1in}

\vspace{0pt}\textbf{Differences in behavior between the force controller and the impedance controller}\par\vspace{0.03in}
{\small The force controller directly commands the desired force, while the impedance controller
 generates force from position and velocity errors. As shown in Figure~\ref{fig:fixed-wall}, 
 the force controller overshoots and oscillates a few times before settling, while the 
 impedance controller is steadier.\par}

\vspace{0.09in}\textbf{Observation and analysis on how the robot behaves under this disturbance}\par\vspace{0.03in}
{\small As shown in Figure~\ref{fig:moving-wall}, the force controller has several large transient spikes but 
then returns near the target force. The impedance controller behaves like a spring-damper system: 
it produces a higher force when the wall moves closer (such as in Figure~\ref{fig:max-amp}) and a lower force when the wall moves 
farther away, and it also has an initial overshoot. This occurs because the impedance controller
uses a virtual spring and damper, which allow it to respond compliantly to the disturbance.\par}

\newpage
\large \textbf{3 - Kinematics}

\vspace{0.1in}
\textbf{End-effector poses for \(q_1\), \(q_2\), and \(q_3\).}
\begin{itemize}
  \item \(q_1\) end-effector pose:\\
  \hspace{0.1in} 
  $$ \textbf{t} = \begin{bmatrix}
    0.088 \\
    0 \\
    0.926
  \end{bmatrix}
  $$
  $$ \textbf{R} = \begin{bmatrix}
    1 & 0 & 0 \\
    0 & -1 & 0 \\
    0 & 0 & -1
  \end{bmatrix}
  $$
  \item \(q_2\) end-effector pose:\\
  \hspace{0.1in} 
  $$ \textbf{t} = \begin{bmatrix}
    0.157 \\
    -0.103 \\
    0.936
  \end{bmatrix}
  $$
  $$ \textbf{R} = \begin{bmatrix}
    0.649 & 0.759 & 0.052 \\
    0.755 & -0.651 & 0.073 \\
    0.089 & -0.008 & -0.996
  \end{bmatrix}
  $$
  \item \(q_3\) end-effector pose:\\
  \hspace{0.1in} 
  $$ \textbf{t} = \begin{bmatrix}
    0.401 \\
    0.087 \\
    0.855
  \end{bmatrix}
  $$
  $$ \textbf{R} = \begin{bmatrix}
    0.980 & -0.181 & -0.081 \\
    -0.174 & -0.593 & -0.786 \\
    0.095 & 0.785 & -0.612
  \end{bmatrix}
  $$
\end{itemize}
\textbf{IK joint angles for (\(R_g, t_g\))}
\begin{itemize}
  \item IK joint angles (rad):\\
  $$ \textbf{q} = \begin{bmatrix}
    0.180637 & 0.257370 & -0.209686 & -2.082814 & -0.032940 & 3.907053 & -3.064487
    \end{bmatrix} 
    $$
\end{itemize}
\textbf{Robot at the IK configuration}


\begin{minipage}[t]{0.495\textwidth}
  \includegraphics[width=\linewidth]{../../HW1/img/Kinematics_IK_pose.png}
  \captionof{figure}{Franka arm with IK joint targets.}
  \label{fig:org-gain}
\end{minipage}\hfill
\begin{minipage}[t]{0.495\textwidth}
  \includegraphics[width=\linewidth]{../../HW1/img/Kinematics_IK_pose_tuned.png}
  \captionof{figure}{Same IK joint targets with tuned gains.}
  \label{fig:new-gain}
\end{minipage}

\vspace{0.1in}
With the original PD gains (\(K_p = 1000, K_d = 1000\)), some joints did not settle at their IK target angles because the actuator torques saturated (Figure~\ref{fig:org-gain}). After reducing the gains (\(K_p = 100, K_d = 20\)), the robot moved closer to the desired joint configuration (Figure~\ref{fig:new-gain}).

% \begin{figure}[htbp]
%   \centering
%   \includegraphics[width=0.5\linewidth]{../../HW1/img/Kinematics_IK_pose.png}
%   \caption{Franka arm at the joint configuration computed by the IK solver.}
% \end{figure}

\end{document}
